\documentclass[10pt]{article}

\usepackage[a4paper,margin=1.15cm,top=0.9cm,bottom=0.8cm,includefoot]{geometry}
\pagestyle{empty}
\usepackage{xcolor}
\usepackage{titlesec}
\renewcommand{\thesection}{\arabic{section}}
\usepackage{enumitem}
\usepackage{hyperref}
\usepackage{fontspec}
\usepackage{xeCJK}

% 中文字体：使用系统已安装的思源黑体（Noto Sans CJK SC）
\setCJKmainfont{Noto Sans CJK SC}[BoldFont={Noto Sans CJK SC Bold}]
\setCJKsansfont{Noto Sans CJK SC}[BoldFont={Noto Sans CJK SC Bold}]
\setmainfont{TeX Gyre Termes}   % 拉丁字母/数字沿用 PDF 常见西文字体

\setlength{\parskip}{0.2em}\setlength{\parindent}{0pt}

\hypersetup{hidelinks}

\definecolor{accent}{HTML}{1A6EBC}
\definecolor{graytext}{HTML}{555555}

\titleformat{\section}{\large\bfseries\color{accent}}{}{0em}{}[{\color{accent}\titlerule[0.6pt]}]
\titlespacing*{\section}{0pt}{0.6em}{0.2em}

\setlist[itemize]{leftmargin=1.3em,itemsep=0.05em,topsep=0.15em,parsep=0pt}

\newcommand{\entry}[3]{%
  \noindent\begin{tabular*}{\textwidth}{@{}l@{\extracolsep{\fill}}r@{}}
    \textbf{#1} & \textcolor{graytext}{#2}
  \end{tabular*}\par\vspace{-0.25em}
  \noindent\textit{#3}\par\vspace{0.22em}
}
\newcommand{\skill}[2]{\textbf{#1}：#2}

\pagestyle{empty}

\begin{document}
\linespread{1.02}\selectfont

{\Huge\bfseries 肖志群}\\[0.3em]
{\color{graytext} 软件工程专业 · 博士研究生}\\[0.5em]
国防科技大学，湖南长沙\\
\href{mailto:xiaozhiqun@nudt.edu.cn}{xiaozhiqun@nudt.edu.cn} \quad$\cdot$\quad
\href{https://zhiqunxiao.github.io}{zhiqunxiao.github.io} \quad$\cdot$\quad
\href{https://github.com/ZhiqunXiao}{github.com/ZhiqunXiao} \quad$\cdot$\quad
\href{https://gitee.com/xiaozq22}{gitee.com/xiaozq22}

\section{教育背景}
\entry{国防科技大学}{2025.09 -- 至今}{软件工程，博士研究生}
\entry{国防科技大学}{2022.09 -- 2025.06}{软件工程，工学硕士}
\entry{中南大学}{2018.09 -- 2022.06}{车辆工程，工学学士}

\section{研究方向}
具身智能（Embodied Intelligence） \quad$\cdot$\quad 智能体机器人（Agentic Robots） \quad$\cdot$\quad 安全人工智能系统（Safe AI Systems）\\
视觉语言动作模型（VLA）\quad$\cdot$\quad 机器人学（Robotics）\quad$\cdot$\quad 任务规划（Task Planning）

\section{学术论文}
\noindent\textbf{已发表}
\begin{itemize}
\item \textbf{OCP: An Online Contingent Planning Method for Robot Tasks with Incomplete Knowledge}\\
  \textit{Z. Xiao}, S. Yang, Y. Xue, S. Wang, X. Mao.\\
  In \textit{Advanced Intelligent Computing Technology and Applications (ICIC 2024)}, Lecture Notes in Computer Science, vol.~14879, pp.~39--50, Springer, 2024, \href{https://doi.org/10.1007/978-981-97-5675-9_4}{doi:10.1007/978-981-97-5675-9\_4}.
\end{itemize}
\vspace{0.2em}
\noindent\textbf{未发表稿件}
\begin{itemize}
\item \textbf{Optimizing Contingent Plans with Parallel Relations and Leveled Plan Graphs} (POCP)\\
  \textit{Z. Xiao}（第一作者），2024。未发表稿件。
\end{itemize}
\vspace{0.2em}
\noindent\textbf{在审稿件}
\begin{itemize}
\item \textbf{一篇面向机器人安全任务规划的研究稿件}（第一作者），2026。目前处于双盲评审阶段，依匿名政策暂不公开题目与投稿去向。
\end{itemize}

\section{实习经历}
\entry{高频交易系统开发实习生}{2022.01 -- 2022.08}{某量化对冲基金，湖南长沙}
\begin{itemize}
\item 负责基础设施组件开发，主管日志系统；
\item 参与交易系统的日常维护。
\end{itemize}

\section{在校经历}
\entry{国防科技大学}{2022.09 -- 至今}{软件工程专业研究生}
\begin{itemize}
\item 2023--2025 年连续三年获校级一等奖学金。
\end{itemize}

\vspace{0.3em}
\entry{中南大学}{2018.09 -- 2022.06}{车辆工程专业本科生}
\begin{itemize}
\item \textbf{大学生方程式赛车队}：担任技术总监助理，团队获 2019 年"蔚来杯"中国大学生电动方程式大赛（FSEC）一等奖。
\item \textbf{大学生创新训练项目（2020）}：主持《基于 FSEC 的四轮独立驱动汽车控制策略》，获国家级立项并顺利结题。
\item \textbf{工程训练综合能力竞赛（2021）}：担任队长，获校二等奖（作品为基于 Arduino 开发的麦克纳姆轮移动机器人）。
\item 2019--2022 年连续三年获校级奖学金。
\item 本科 GPA：3.7/4.0（专业前 30\%）。
\end{itemize}

\enlargethispage{2\baselineskip}
\section{专业技能与兴趣爱好}
\begin{itemize}
\item \skill{编程语言}{C/C++、Python、MATLAB、Simulink}
\item \skill{嵌入式开发}{Arduino、树莓派（Raspberry Pi）}
\item \skill{系统与开发工具}{Linux、Shell、Git、Docker}
\item \skill{机器学习与智能体框架}{PyTorch、Scikit-learn、LangChain、LangGraph}
\item \skill{机器人系统}{ROS/ROS2、Isaac Sim、MuJoCo；具备 TurtleBot3、TIAGo 等多款移动机器人开发经验}
\item \skill{机械设计}{Catia、SolidWorks、AutoCAD}
\item \skill{个人兴趣}{徒步、乒乓球、游泳}
\end{itemize}

\end{document}
